\section{数学分析课程内不同知识点的（相互）关联}
\subsection*{话题 1-A：积分与微分的对应}
例 7.1.1 的积分不等式与例 7.1.2 的微分不等式可由 Newton--Leibniz 公式互相转化：前者的 $f$ 对应后者的 $f'$。

\ex{7.1.1}设 $f\in C^1[a,b]$，$f(a)=f(b)=0$。证明
\[\max_{a\le x\le b}|f'(x)|\ge \frac{4}{(b-a)^2}\left|\int_a^b f(x)\dd x\right|.\]
\pf 由两端点积分得 $|f(x)|\le M\min\{x-a,b-x\}$，其中 $M=\max|f'|$。积分并利用 $\int_a^b\min\{x-a,b-x\}\dd x=(b-a)^2/4$ 即得。

\ex{7.1.2}设 $f\in C^2[a,b]$，$f'(a)=f'(b)=0$。证明
\[\max_{a\le x\le b}|f''(x)|\ge\frac4{(b-a)^2}|f(b)-f(a)|.\]
\pf 对 $f'$ 应用例 7.1.1。

\ex{7.1.3}设 $f\in C[a,b]$，在 $(a,b)$ 可微，且 $f(a)=f(b)=0$。证明存在 $\xi\in(a,b)$，使
\[|f'(\xi)|=\frac4{(b-a)^2}\left|\int_a^b f(x)\dd x\right|.\]
\pf 令右端为 $c$。若导数的绝对值处处不等于 $c$，利用导数的介值性与端点条件，可把 $|f'|$ 全局限制于 $c$ 的一侧；再用例 7.1.1 的积分估计或 Rolle 定理得到矛盾。

\ex{7.1.4}设 $f$ 在 $[a,b]$ 可导，在 $(a,b)$ 二阶可导，$f'(a)=f'(b)=0$。证明存在 $\xi\in(a,b)$，使
\[|f''(\xi)|=\frac4{(b-a)^2}|f(b)-f(a)|.\]
\pf 记右端为 $k$。用二阶导数的介值性讨论 $f''$ 与 $\pm k$ 的关系。若在 $(a,b)$ 上始终有 $|f''|<k$，则由端点导数为零推出 $f'$ 在闭区间连续，转化为例 7.1.3。另一种方法是选 $\beta$ 使 $f(\beta)=(f(a)+f(b))/2$，分别在 $[a,\beta]$ 和 $[\beta,b]$ 使用二阶中值公式，得到
\[f''(\xi_1)=\frac{f(b)-f(a)}{(\beta-a)^2},\qquad f''(\xi_2)=-\frac{f(b)-f(a)}{(b-\beta)^2}.\]
二者绝对值分别落在 $k$ 的两侧，再由介值性得到结论。

\ex{7.1.5}在例 7.1.2 的条件下，若 $f(a)\ne f(b)$，则严格有
\[\max_{[a,b]}|f''|>\frac4{(b-a)^2}|f(b)-f(a)|.\]
\ex{7.1.6}在例 7.1.4 的条件下，若 $f(a)\ne f(b)$，则存在 $\xi\in(a,b)$，使
\[|f''(\xi)|>\frac4{(b-a)^2}|f(b)-f(a)|.\]
\pf 若等号可达到最优估计，则 $f'$ 必为以区间中点为顶点的折线型函数；它不能同时满足所需的二阶可导性。结合例 7.1.4 和导数介值性得到严格不等式。
\par 这组例子展示了如何从已知的积分估计生成微分题。例 7.1.3 比例 7.1.1 的结论强；再对 $f'$ 作对应替换得到例 7.1.4。困难在于 $f'$ 不一定在闭区间连续，因此不能机械地套用积分公式。二阶导数的 Darboux 性质在这里补上了连续性缺口。若只要得到例 7.1.2，可在例 7.1.4 的二阶中值公式中直接取 $\beta=(a+b)/2$。

\subsection*{话题 1-B：导数与差分估计}
\ex{7.1.7}有界数列 $\{a_n\}$ 满足 $a_{n+2}-2a_{n+1}+a_n\to0$。证明 $a_{n+1}-a_n\to0$。
\pf 固定整数 $k\ge2$，将 $a_{n+k}-a_n$ 用首个差分与二阶差分展开，可得
\[k|\Delta a_n|\le 2\sup_{j\ge n}|a_j|+\frac{k(k-1)}2\sup_{j\ge n}|\Delta^2a_j|.\]
先令 $n\to\infty$，再令 $k\to\infty$。

\ex{7.1.8}若 $f$ 在 $[0,\infty)$ 有界、二阶可导，且 $f''(x)\to0$，则 $f'(x)\to0$。
\pf 固定 $h>0$，Taylor 中值估计给出 $|f'(x)|\le 2\|f\|_\infty/h+(h/2)\sup_{t\ge x}|f''(t)|$。先让 $x\to\infty$，再让 $h\to\infty$。

\ex{7.1.9}若 $f$ 在 $[a,\infty)$ 二阶可导，记 $M_j=\sup_{x\ge a}|f^{(j)}(x)|$（$j=0,1,2$），则存在与 $f$ 无关的常数 $C$，使 $M_1^2\le CM_0M_2$（允许部分 $M_j=+\infty$，但 $\{M_0,M_2\}\ne\{0,+\infty\}$）。
\pf 对 $x\ge a$ 和 $h>0$，Taylor 公式给出
\[|f'(x)|\le \frac{|f(x+h)-f(x)|}{h}+\frac h2M_2\le\frac{2M_0}{h}+\frac h2M_2.\]
取 $h=2\sqrt{M_0/M_2}$，得 $M_1\le2\sqrt{M_0M_2}$。端点退化情形由极限处理。
\ex{7.1.10}若 $f$ 在 $[0,\infty)$ 二阶可导、$f''$ 有界且 $f(x)\to0$，则 $f'(x)\to0$。
\ex{7.1.11}若 $f,f''$ 在 $[a,\infty)$ 有界，则 $M_1^2\le4M_0M_2$。
\ex{7.1.12}若 $f,f''$ 在 $\R$ 上有界，则 $M_1^2\le2M_0M_2$。
\pf 由 Taylor 公式分别对单侧或双侧的 $f(x\pm h)$ 作估计，优化 $h$ 即得。

\ex{7.1.13}记 $\Delta a_n=a_{n+1}-a_n$，$\Delta^2a_n=a_{n+2}-2a_{n+1}+a_n$。证明存在常数 $C$，对任意整数 $k\ge2$ 有
\[k\sup_n|\Delta a_n|\le C\left(\sup_n|a_n|+k^2\sup_n|\Delta^2a_n|\right).\]
等价地，对任意 $\alpha>0$，
\[\sup_n|\Delta a_n|\le C\left(\alpha\sup_n|a_n|+\frac1\alpha\sup_n|\Delta^2a_n|\right),\]
并可推出 $\sup|\Delta a_n|^2\le C\sup|a_n|\sup|\Delta^2a_n|$。

\ex{7.1.15}若 $f$ 在 $[a,\infty)$ 有 $n+1$ 阶导数，$M_0,M_{n+1}<\infty$，则对 $1\le k\le n$ 有
\[M_k\le C_k M_0^{1-k/(n+1)}M_{n+1}^{k/(n+1)}.\]
\ex{7.1.16}定义 $\Delta^{k+1}a_n=\Delta^ka_{n+1}-\Delta^ka_n$。对 $0<j<k$，存在常数 $C_{k,j}$ 使
\[\alpha^j\sup_n|\Delta^ja_n|\le C_{k,j}\left(\sup_n|a_n|+\alpha^k\sup_n|\Delta^ka_n|\right),\quad\alpha>0.\]
\par 例 7.1.15 和例 7.1.16 分别给出导数与差分的高阶插值估计。调整尺度 $\alpha$ 可以将两个端点阶数的界转为中间阶数的界。这也解释了例 7.1.7 与 7.1.8 的相似性：有界函数的高阶变化越来越小，就会迫使低阶变化趋零。

\ex{7.1.17}设 $f$ 在 $[0,\infty)$ 有界、三阶可导，$f'''(x)\to\alpha\in\R$。则 $f''(x)\to0$，进而 $\alpha=0$。
\ex{7.1.18}若 $f(x)/x\to0$ 且 $x|f''(x)|$ 在无穷远处有界，则 $f'(x)\to0$。
\pf 在 $[x,(1+\delta)x]$ 上应用中值定理，得 $|f'(x)|\le |f((1+\delta)x)-f(x)|/(\delta x)+O(\delta)$；先令 $x\to\infty$，再令 $\delta\downarrow0$。
\ex{7.1.20}若 $0<\alpha<1$，$f(x)=o(x^\alpha)$ 且 $x^\alpha|f''(x)|$ 有界，证明 $f'(x)\to0$；并讨论 $\alpha=1$ 时的结论。
\ex{7.1.21}若 $0<\alpha<1$，$f(x)=O(x^\alpha)$ 且 $x^\alpha f''(x)\to0$，证明 $f'(x)\to0$；并讨论 $\alpha=1$。
\ex{7.1.22}若 $f(x)\to0$，且 $f''+\lambda f'$ 在 $[0,\infty)$ 上有界（$\lambda$ 为常数），证明 $f'(x)\to0$。
\pf 把 $g=f''+\lambda f'$ 看作已知有界函数，对一阶方程 $f''+\lambda f'=g$ 作积分因子估计，再结合 $f\to0$ 与差分商估计即可。
\par 对这些变式，应注意 $f(x)\to0$ 与 $f''$ 的全局有界是不同层面的条件。对于 $x^\alpha$ 权重，常可在长度为 $\delta x$ 的区间上使用差分商；二阶余项为 $O(\delta x\sup|f''|)$，让 $x$ 先趋于无穷，再让 $\delta$ 趋于零。

\ex{7.1.23}设 $p\in[1,\infty)$，$f\in C^2[a,b]$。证明存在 $C_p$ 使
\[\|f'\|_{L^p}\le C_p(\|f\|_{L^p}+\|f''\|_{L^p}),\qquad \|f'\|_{L^p}^p\le C_p(\|f\|_{L^p}^p+\|f''\|_{L^p}^p).\]
\ex{7.1.24}在例 7.1.23 的条件下，对任意 $\varepsilon>0$，存在 $C_{p,\varepsilon}$ 使
\[\|f'\|_{L^p}\le C_{p,\varepsilon}\|f\|_{L^p}+\varepsilon\|f''\|_{L^p},\quad
\|f'\|_{L^p}^p\le C_{p,\varepsilon}\|f\|_{L^p}^p+\varepsilon\|f''\|_{L^p}^p.\]
\ex{7.1.25}设 $f\in C^2[0,1]$ 且 $f(0)f(1)\ge0$。证明
\[\int_0^1|f'(x)|\dd x\le2\int_0^1|f(x)|\dd x+\int_0^1|f''(x)|\dd x.\]
\par 上述三个积分型不等式说明，不能只依靠形式上的 $\|f'\|_p\le C(\|f\|_p+\|f''\|_p)$。题目给定的端点符号条件会影响最佳常数；在应用基本定理时，需要把区间分成单调段并结合边界值。

\subsection*{话题 1-C：单调有界的矩阵列}
\ex{7.1.26}若 $n$ 阶实矩阵列 $A_k$ 按半正定序单调递增，且存在矩阵 $M$ 使 $A_k\preceq M$，则 $A_k$ 收敛。
\pf 对任意 $x\in\R^n$，标量列 $x^\mathsf TA_kx$ 单调有界而收敛。由极化恒等式可知各对称矩阵元均收敛。若未事先要求 $A_k$ 自身对称，则 $A_{k+1}-A_k$ 为半正定矩阵，故反对称部分不随 $k$ 变化；全部矩阵元仍收敛。
\ex{7.1.27}设 $\{P_k\}$ 为 $[a,b]$ 上单调递增、一致有界且次数不超过固定整数 $n$ 的实多项式列。则 $P_k$ 一致收敛到某个次数不超过 $n$ 的多项式。
\pf 取 $n+1$ 个互异点。各点处的函数值单调有界而收敛；插值多项式的系数随这 $n+1$ 个值连续变化，故系数逐项收敛。有限维空间中系数收敛推出一致收敛。
\ex{7.1.28}把例 7.1.27 推广到 $\R^m$ 区域 $\Omega$ 上的次数不超过 $n$ 的多项式列：若逐点单调递增且局部一致有界，则局部一致收敛到次数不超过 $n$ 的多项式。
\pf 在 $\Omega$ 内选择足够多的插值点，使多元多项式的系数由这些点的值唯一确定；其余同例 7.1.27。
\ex{7.1.31}把例 7.1.27 的“多项式”换成次数不超过 $n$ 的 $2\pi$ 周期三角多项式，仍可得到同样的结论。
\pf 所讨论的三角多项式组成有限维线性空间；选取足够多的点值作为坐标，再使用有限维性。
\ex{7.1.33}设 $f_k$ 是 $[a,b]$ 上连续、逐点单调递增且一致有界的凸函数列。证明 $f_k$ 在 $[a,b]$ 上一致收敛到凸函数。
\pf 逐点单调有界给出极限 $f$，凸性在极限下保留。有限凸函数在内部连续。端点处，单调递增的连续函数列的极限为下半连续函数；而凸函数在闭区间端点可能出现的间断只能是向上的跳跃，该跳跃与下半连续性矛盾。所以 $f$ 在闭区间连续。Dini 定理推出一致收敛。
